\begin{proposition}\label{prop:decic-no78}
An elliptic curve over a number field of degree $10$ has no point
of order $78$, and its torsion does not contain $C_3\times C_{39}$.
\end{proposition}
\begin{proof}
Put $X=X_1(39)$, of genus $33$.
Its rational Jacobian is finite, and specialization at $2$ is
injective on it, as recalled in Proposition~\ref{prop:decic-no117}.
Its rational gonality is $14$.
We determine the noncuspidal rational effective divisors of degree
$10$ on $X$.
Let $D$ be such a divisor. The global and finite cusp signatures
are both $1^{12}2^6 12^2$. Hence every cuspidal divisor of degree
at most $10$ in the special fibre lifts globally.
An entirely cuspidal reduction would give $D\sim C$ for its
rational cuspidal lift $C$, by injectivity. Gonality $14>10$
would then give $D=C$, a contradiction.

The complete noncuspidal place counts over $\mathbb F_2$, in
degrees $1,\ldots,10$, are
\[
                       (0,0,0,0,0,12,0,49,52,60).
\]
Thus $\overline D=P+C$, with $P$ a noncuspidal place of degree
$6,8,9$ or $10$ and $C$ a cuspidal divisor of the complementary
degree. Every such $C$ lifts globally.
Put $A_5=T_5-1-5\langle5\rangle$, using point levels.
It annihilates $J_X(\mathbb Q)$ by Eichler--Shimura and reduction
at $5$, and it annihilates rational cusp divisors by the cusp
boundary formula. Consequently $A_5P$ must be principal.
The accepted degree-independent tests exclude the sextic orbit
and four of the five octic orbits. The remaining octic orbit is
the reduction of the known closed divisor $D_8$.
Injectivity and gonality would give $D=hD_8+C$, contradicting
noncuspidal support.

There are five diamond orbits in each of degrees $9$ and $10$.
On the explicit degree-$17$ function-field model of genus $33$,
the exact computations give
\[
\begin{array}{c|ccccc}
\deg P&\multicolumn{5}{c}{h^0(A_5P)\text{ on the five source orbits}}\\
\hline
9&1&0&0&0&0\\
10&0&0&0&0&1.
\end{array}
\]
The degree-$9$ survivor is precisely the diamond orbit of the
reduction of the known closed nonic divisor $D_9$.
The degree-$10$ survivor is precisely the orbit of the reduction
of the known decic divisor $D_{10}$ represented by the point
$\texttt{1.39.10.a}$.
Both identifications use exact finite-field arithmetic on the
known characteristic-zero equations; their defining polynomials
remain irreducible modulo $2$.
A fresh native reconstruction verifies all ten Riemann--Roch
spaces and all $49$ prime references using two-generator ideal
identities. Independent exhaustive point enumeration checks the
$52$ and $60$ places and their orbit coverage.
Independently constructed marked Velu maps verify all $60$
geometric $5$-kernels, their source and diamond coordinates,
and every residue norm multiplicity.

If $\deg P=9$, injectivity and gonality give $D=hD_9+C$,
again contradicting noncuspidal support.
If $\deg P=10$, the same argument gives $D=hD_{10}$.
Thus every noncuspidal rational effective divisor of degree $10$
on $X$ is a diamond translate of $D_{10}$.

Suppose now that $E/K$ has a point $Q$ of order $78$, with
$[K:\mathbb Q]=10$, and apply this classification to the trace
of $(E,2Q)$. It equals a closed divisor $hD_{10}$, so the moduli
field has degree $10$ and is $K$.
The fine moduli interpretation identifies $E$ with a conjugate
of the known decic curve, over the corresponding residue field.
That curve has full torsion $C_{39}$: its exact order-$39$ point
and good unramified odd-prime reductions with torsion-order
bound $39$ are independently verified.
This contradicts the point $Q$ of order $78$.
The same classification excludes $C_3\times C_{39}$, since the
known decic curve has no independent second point of order $3$.
\end{proof}
